% !TEX root = ./NIR.tex
Проведение современных астрономических наблюдений предполагает высокую точность наведения и сопровождения телескопов, ошибки в определении углового положения элементов системы приводят к смещению изображения небесного объекта и снижению качества получаемых данных. Для обеспечения стабильного и высокоточного позиционирования применяются специальные датчики углового положения, преобразующие пространственную ориентацию в цифровой сигнал. Среди них наибольшее распространение получили \emph{абсолютные энкодеры}, обеспечивающие однозначное определение углового положения без необходимости повторной калибровки при прерывании работы. Важное преимущество абсолютного энкодера это отсутствие необходимости инициализации путем поворота до срабатывания нулевой точки отсчета, что требует дополнительного времени и, в общем случае, небезопасно.

Актуальность разработки и совершенствования таких устройств обусловлена постоянным ростом требований к точности, стабильности и надежности навигационных систем астрономических установок. Современные телескопы, в зависимости от задачи, требуют измерения углов с погрешностью от единиц до десятых долей угловых секунд при сохранении устойчивости к внешним воздействиям --- вибрациям, температурным колебаниям и загрязнению. При этом конструкция датчика должна быть компактной и обеспечивать возможность интеграции в существующие механические узлы.

Большинство промышленных решений в области абсолютного позиционирования разработаны для применения в робототехнике, станках и системах автоматизации, где требования к точности и термостабильности менее жесткие, чем в астрономии. Использование таких систем в телескопах зачастую требует дополнительной адаптации, а стоимость специализированных энкодеров известных производителей (Heidenhain, Renishaw, Celera Motion и др.) остается высокой. Поэтому актуальной задачей является исследование принципов построения и создание прототипа абсолютного датчика угла, оптимизированного для условий эксплуатации астрономических установок.  

Целью данной работы является разработка и исследование прототипа абсолютного датчика угла для астрономического телескопа на основе оптического принципа измерения.

Задачи работы:
\begin{itemize}
    \item провести анализ существующих технологий и конструкций абсолютных энкодеров, применяемых в высокоточных системах;  
    \item определить ключевые метрологические и эксплуатационные требования к датчику угла для астрономических установок;  
    \item разработать концепцию и изготовить прототип абсолютного энкодера; 
    \item провести экспериментальные исследования полученного образца и оценить его характеристики.  
\end{itemize}

Таким образом, структура курсовой работы включает теоретическую часть, содержащую обзор и анализ существующих технологий, и практическую часть, посвященную разработке и исследованию прототипа. Итогом работы станет реализация оптического абсолютного датчика угла с использованием доступных компонентов при сохранении необходимого уровня точности для применения в телескопных системах.  